\documentclass[10pt,a4paper]{amsart}
\usepackage[margin=19mm]{geometry}
\usepackage{amsmath,amssymb,mathtools}
\usepackage[scr=rsfs,cal=euler]{mathalfa}
\usepackage{xfrac}
\usepackage[unicode,hidelinks,bookmarks=false]{hyperref}
\newcommand{\ie}{i.e.\ }
\newcommand{\cf}{cf.\ }
\newcommand{\resp}{respectively\ }
\newcommand{\Subsection}{\subsection}
\providecommand{\nobreakdash}{\nobreak\hbox{-}}
\setlength{\emergencystretch}{2em}
\multlinegap0pt
\usepackage{amssymb}
\usepackage{mathtools}
\newtheorem{lemma}{Lemma}
\newtheorem{proposition}{Proposition}
\title{Rings and Boolean Algebras as Algebraic Theories}
\author{Arturo De Faveri}
\date{Source: arXiv:2503.04430v3 (CC BY 4.0)}
\begin{document}
\maketitle

\noindent\emph{Source and licence.} Rings and Boolean Algebras as Algebraic Theories. Version arXiv:2503.04430v3 (CC BY 4.0), distributed by arXiv under the Creative Commons Attribution 4.0 International licence, http://creativecommons.org/licenses/by/4.0/, as recorded in the arXiv licence metadata for this exact version and shown on the abstract page https://arxiv.org/abs/2503.04430v3. Full licence notice: \texttt{licenses/arxiv-2503.04430-CC-BY-4.0.txt}.\par\smallskip

\noindent\textit{Editorial scope.} Counted unit: the article's proposition prop:coord (source0.tex lines 926--936). The article's complete original proof of it is delivered with it (source0.tex lines 937--989, one \texttt{proof} environment of 53 lines). No mathematics is written, repaired or re-derived: the statement and the proof are the article's own lines, in the article's own notation, and no retained line is substituted or re-flowed.  Retained uncounted as the source-local context the proof needs: lines 854--860.  Nothing was omitted from any retained span, so the package carries no omission record.  No retained line contains a citation, so the wrapper carries no bibliography and no bibliography entry.  No retained line contains a figure environment, an image-inclusion command or drawing code, so no image is delivered and the package declares no asset dependency.  Editorial formatting only, in addition: an AMS \texttt{amsart} preamble that restates the article's own declarations and macros used by the retained lines, namely the environments \texttt{lemma}, \texttt{proposition} on separate counters, together with the standard packages those lines need.  The retained environments print in this document as Lemma 1, Proposition 1 in order of appearance, whereas the article prints the same items with the numbering of its own full document.  The complete native source of the article as distributed in the arXiv e-print for this exact version is bundled unchanged as \texttt{sources/174328/source0.tex}.
\begin{lemma}
  \label{lem:pl}
  If $T$ is an affine theory, then 
  \begin{equation*}
    p(a(x,y), b(x,y), c(x,y))=(a-b+c)(x,y)\text{.}
  \end{equation*}
\end{lemma}

\begin{proposition}
  \label{prop:coord}
  Let $T$ be an affine theory.
  Let $R$ be the commutative ring $T(2)$. 
  The following hold: 
  \begin{enumerate}
    \item if $f \in T(n)$, then $f[1]+\cdots+f[n]=1$ in $R$; 
    \item for $f,g \in T(n)$, if $f[i]=g[i]$ for all $1 \le i \le n$, then $f=g$; 
    \item for every $(r_1, \ldots, r_n) \in R^n$ such that $r_1+\cdots+r_n=1$, there is $f \in T(n)$ such that $f[i]=r_i$ for all $1 \le i \le n$. 
  \end{enumerate}
\end{proposition}

\begin{proof}
    Regarding the first item, we prove that $f[1] + \cdots + f[n] = 1$. 
    Observe that 
    \begin{align*}
      (f[1]+f[2])(x,y) & = p(f[1](x,y),y,f[2](x,y)) & \\
      & = p(f(x,y, \ldots, y) , f(y,y \ldots, y), f(y,x, \ldots, y)) & f \text{ is idemp.}\\
      & = f(p(x,y,y), p(y,y,x), \ldots, p(y,y,y)) & p,f \text{ commute} \\
      & = f(x,x,y \ldots, y) & p \text{ is Mal'cev} 
    \end{align*}
    and, therefore, inductively, 
    \begin{equation*}
      (f[1]+ \cdots + f[n])(x,y) = f(x,x, \ldots,x) =x \text{.}
    \end{equation*}

    For the second item, observe that 
    \begin{align*}
      p(f[1](x_1,y),y,f[2](x_2,y)) & = p(f(x_1,y, \ldots, y) , f(y,y \ldots, y), f(y,x_2, \ldots, y)) & \\
      & = f(p(x_1,y,y), p(y,y,x_2), \ldots, p(y,y,y)) &  \\
      & = f(x_1,x_2,y \ldots, y) &
    \end{align*}
    and, therefore, inductively, 
    \begin{equation*}
      p(f(x_1,\ldots,x_{n-1},y), y,f[n](x_n,y)) = f(x_1,x_2, \ldots,x_n)\text{.}
    \end{equation*}
    This implies that if $f[i]=g[i]$ for all $i=1, \ldots, n$, then $f = g$. 

  Finally, we prove the third item. 
  %First of all, observe that, if $(r_1,r_2,r_3)=(1,-1,1)$, then any Mal'cev operation $p \in T(3)$ is such that $p[i]=r_i$ for $i=1,2,3$.  
  The proof is by induction on $n$. 
  If $n=1$, then clearly $f(x)=x$. 
  Assume that the result holds for $n-1$, and let $(r_1, \ldots, r_n) \in R^n$ be such that $r_1 + \cdots +r_n=1$. 
  By inductive assumption there is $g \in T(n-1)$ such that $g[1]=r_1 + r_2$ and $g[i]=r_{i+1}$ for $2 \le i \le n$.
  Let $f(x_1, \ldots, x_n):=p(r_2(x_2,x_1),x_1,g(x_1,x_3,\ldots,x_n))$. 
  We prove that $f[i]=r_i$; we separate three cases $i=1,2$ and $3 \le i \le n$:  
  \begin{align*}
    f(x,y,\ldots,y) & = p(r_2(y,x),x,g(x,y,\ldots,y)) & \\
    & = r_2(y,x) -x +g[1](x,y) & \text{ by Lemma \ref{lem:pl}}\\
    & = (1-r_2 -1+r_1+r_2)(x,y) & \text{ by inductive assumption} \\
    & = r_1(x,y)
  \end{align*} 
  \begin{align*} 
    f(y,x,y,\ldots,y) & = p(r_2(x,y),y,g(y, \ldots, y)) & \\
    & = p(r_2(x,y),y,y) & g \text{ is idempotent}\\
    & = r_2(x,y) & p \text{ is Mal'cev} 
  \end{align*}
  \begin{align*}
    f[i](x,y) & = p(r_2(y,y),y,g[i-1](x,y)) & 3 \le i \le n \\
    & = p(y,y,g[i-1](x,y)) & r_2 \text{ is idempotent} \\
    & = g[i-1](x,y) & p \text{ is Mal'cev} \\
    & = r_i(x,y) & \text{ by inductive assumption} 
  \end{align*}
  This concludes the proof. 
\end{proof}

\end{document}
